\section*{Заключение}
\addcontentsline{toc}{section}{Заключение}
 
В данной работе проведён сравнительный анализ четырёх методов статистического
отбора~--- \textit{случайного, механического, типического и кластерного}~--- применительно
к задаче оценки потребительских предпочтений в отношении аудио-брендов
Sony, JBL и Marshall по модели идеальной точки. По результатам исследования
можно сформулировать следующие выводы.

\underline{Лидер рынка однозначно определeн}. По генеральной совокупности
бренд Sony имеет наименьшее среднее отклонение от идеальной точки. Этот
вывод устойчив: во всех семи конфигурациях имитационного моделирования
и при $10\,000$ итераций именно Sony занимала первое место по пользовательским
предпочтениям.
 
Объема разведочной выборки в $4\%$ от генеральной совокупности \\
\underline{недостаточно для разграничения лидеров}.
Sony и JBL статистически
неразличимы~--- они делят победы примерно поровну по всем методам.
Это подтверждает необходимость применения итерационной схемы расчёта
минимально необходимого объёма выборки $n_{min}$, после достижения которого
неопределённость устраняется.
 
Можно заключить, что для маркетинговых исследований в условиях известной
поло-возрастной структуры аудитории рекомендуется применять
\underline{типический повторный отбор}: он контролирует состав выборки,
исключает межгрупповую составляющую вариации и обеспечивает наиболее
стабильные и несмещённые оценки интегрального показателя брендов.
 
Итогом работы стала реализация автоматизированного цикла обработки данных, интегрирующего 
методы математической статистики с алгоритмами бизнес-аналитики. Созданный программный комплекс
позволяет не только рассчитывать рейтинги брендов, но и количественно оценивать надежность 
полученных  результатов, обеспечивая воспроизводимость исследованиях
на произвольных структурах данных
